\section{Method}
\label{sec:method}

\begin{figure*}[t]
\centering
\resizebox{0.9\textwidth}{!}{%
\begin{tikzpicture}[
  font=\small,
  box/.style={draw, rounded corners=2pt, minimum height=0.9cm, align=center, inner sep=4pt},
  fz/.style={box, fill=frozen!15, draw=frozen, thick},
  tr/.style={box, fill=trained!15, draw=trained, thick},
  tok/.style={draw, minimum width=0.38cm, minimum height=0.42cm, inner sep=0pt},
  arr/.style={-{Stealth[length=2.2mm]}, thick}]
  % inputs and the two streams
  \node[box, fill=gray!10, minimum width=1.9cm] (x) at (0,0) {Meeting\\audio $\bm{x}$};
  \node[fz, minimum width=2.7cm] (enc) at (3.2,1.0) {Audio encoder\\(Qwen3-Omni)};
  \node[fz, minimum width=2.7cm] (dia) at (3.2,-1.0) {Simultaneous diarizer\\(Nemotron 3), $S$ slots};
  \node[tr, minimum width=2.3cm] (proj) at (6.5,-1.0) {Speaker projector\\$g_\psi$ (\S\ref{sec:method-thinker})};
  \draw[arr] (x.east) -- ++(0.3,0) |- (enc.west);
  \draw[arr] (x.east) -- ++(0.3,0) |- (dia.west);
  \draw[arr] (dia) -- node[above]{$\bm{P}$} (proj);
  % interleaved sequence
  \foreach \i/\c in {0/frozen,1/frozen,2/frozen,3/frozen,4/frozen,5/spkc,6/frozen,7/frozen,8/frozen,9/frozen,10/frozen,11/spkc} {
    \node[tok, fill=\c!35] (t\i) at (8.9+\i*0.4,0) {};
  }
  \node (dots) at (8.9+12*0.4+0.05,0) {$\cdots$};
  \draw[decorate, decoration={brace, amplitude=4pt}] (t0.north west) -- node[above=4pt, font=\small]{$\bm{A}_1$: $K$ acoustic tokens} (t4.north east);
  \node[below=2pt of t5, font=\small] {$\bm{s}_1$};
  \node[below=2pt of t11, font=\small] {$\bm{s}_2$};
  \node[font=\small] at (11.2,-1.25) {interleaved input $\bm{E}$ (Eq.~\ref{eq:interleave})};
  \draw[arr] (enc.east) -| node[pos=0.3, above, font=\small]{acoustic tokens $\bm{a}_i$} (t8.north);
  \draw[arr] (proj.east) -- node[above, font=\small]{$\bm{s}_j$} ++(0.75,0) |- (t0.west);
  % LLM
  \node[tr, minimum height=1.9cm, minimum width=2.0cm] (llm) at (15.6,0) {Thinker LLM\\+ LoRA\\(\S\ref{sec:method-thinker})};
  \draw[arr] (dots.east) -- (llm.west);
  % output segments (second row)
  \node[box, fill=white, anchor=west] (o1) at (6.2,-2.6)
    {$\spk{2}\,\tim{5}$\ \zh{我想知道会议的时间}\ $\tim{48}$};
  \node[box, fill=white, anchor=west] (o2) at (6.2,-3.7)
    {$\spk{3}\,\tim{42}$\ \zh{对，没错。}\ $\tim{56}$};
  \draw[arr] (llm.south) -- ++(0,-0.85) -| ([xshift=-0.3cm]o1.west) -- (o1.west);
  \draw[arr] ([xshift=-0.3cm]o1.west) |- (o2.west);
  \node[font=\small, anchor=west] at (6.2,-4.45) {who, when, what (Eq.~\ref{eq:serial})};
  % talker
  \node[tr, minimum width=2.4cm] (talk) at (14.9,-3.15) {PA Talker, voice $\bm{v}_k$\\one stream per speaker\\(\S\ref{sec:method-queue}--\ref{sec:method-train})};
  \draw[arr] (o1.east) -- ++(0.25,0) |- (talk.west);
  \draw[arr] (o2.east -| o1.east) ++(0,0) -- ++(0.25,0) |- (talk.west);
  \draw[thick] (o2.east) -- (o2.east -| o1.east);
  \node[box, fill=gray!10, minimum width=1.7cm] (y) at (17.6,-3.15) {Translated\\speech};
  \draw[arr] (talk) -- (y);
  % voice prompts
  \draw[arr, dashed, spkc!80!black] (dia.south) -- ++(0,-3.75) -| node[pos=0.25, below, font=\small]{voice prompt $\bm{v}_k$: solo regions of slot $k$ (\S\ref{sec:method-queue})} (talk.south);
  % simultaneous read loop
  \draw[arr, dotted, trained!80!black] (llm.north) -- ++(0,0.5) -- node[above, font=\small]{simultaneous: read every $\Delta$\,s, write or $\wait$ (\S\ref{sec:method-thinker})} ++(-7.4,0);
\end{tikzpicture}}
\caption{Overview of \sys{}. Acoustic tokens (what is said) are interleaved with speaker tokens that a trained projector computes from the slot activities $\bm{P}$ of a frozen simultaneous diarizer (who says it). Reading every $\Delta$ seconds, the Thinker writes speaker-first, time-stamped segments; the PA Talker speaks them in one stream per speaker and voice, trained on the posterior of the playback-alignment CRF, which does not run at test time. Blue: frozen; orange: trained.}
\label{fig:arch}
\end{figure*}


For every contribution to a meeting, \sys{} decides who said what (\S\ref{sec:method-thinker}) and when the listener
hears its translation (\S\ref{sec:method-queue}--\ref{sec:method-train}); Figure~\ref{fig:arch} gives an overview.

\subsection{Problem Formulation}
\label{sec:method-task}

A translation helps a meeting participant only if it is attributed to the right speaker and played soon after that
speaker spoke. The input is a single-channel recording $\bm{x}$ of $T$ seconds with an unknown number of, sometimes
overlapping, speakers. The output is a sequence of segments $\big((k_n,a_n,b_n,\bm{y}^{(n)})\big)_{n=1}^{N_\mathcal{Y}}$:
a speaker label $k_n\in\{1,\dots,S\}$ that stays with one person for the whole meeting, the start and end $a_n<b_n$ of
the source speech, and its translation $\bm{y}^{(n)}$. The system reads the recording at times $t_1<\dots<t_R=T$; what it
commits at read $r$ depends only on $\bm{x}_{\le t_r}$ and is never revised. Each translation is spoken in its
speaker's voice, in a stream of that speaker alone; the time at which a target unit (word or character) is played
there is its \emph{heard time}, computed from the playback schedule (\S\ref{sec:method-queue}), not measured on
listeners. A translation is thus delayed twice: until it is committed, and from then until it is heard.

\subsection{Speaker-Attributed Simultaneous Translation}
\label{sec:method-thinker}

The acoustic tokens of a speech LLM describe what is said but not who says it: they keep no identity across a meeting,
and concurrent speakers share them. A translator that sees only these tokens can merge concurrent contributions, and a
contribution it never translates cannot be attributed afterwards (Figure~\ref{fig:motivation}). \sys{} therefore gives
the Thinker a meeting-wide speaker stream and lets it name the speaker before translating.

\paragraph{Speaker tokens.} The frozen audio encoder of Qwen3-Omni~\citep{xu2025qwen3omni} maps the recording to
acoustic tokens $\bm{A}\in\R^{N\times d}$, where $d$ is the hidden size of its LLM, the Thinker. A frozen simultaneous
diarizer, Nemotron 3 Diarization~\citep{medennikov2026nemotron}, returns speaker-activity posteriors
$\bm{P}\in[0,1]^{F\times S}$ over $F$ frames and $S$ slots; slots are ordered by first arrival and kept for the whole
meeting, and overlapped speech activates several of them. A projector $g_\psi$ maps the activities
$\bm{P}_j\in[0,1]^{K\times S}$ resampled to the $j$-th block of $K$ acoustic tokens to a speaker token
$\bm{s}_j=g_\psi(\bm{P}_j)\in\R^{d}$, which, following G-STAR~\citep{peng2026gstar}, is placed after the block it describes:
\begin{equation}
\bm{E}=[\bm{A}_1;\bm{s}_1;\dots;\bm{A}_M;\bm{s}_M]\in\R^{(N+M)\times d},
\label{eq:interleave}
\end{equation}
with $\bm{A}_j\in\R^{K\times d}$ and $M=\lceil N/K\rceil$, and fills the audio slot of the Thinker's prompt
(Appendix~\ref{app:train-thinker}).

\paragraph{Speaker-first, time-stamped output.} The output must then say who is speaking. We add a speaker token
$\spk{k}$ per slot and time tokens $\tim{m}$, marking $m\delta_t$ seconds from the window start, and serialise a
window's segments by start time, speaker first, as in serialised output training~\citep{kanda2020sot}:
\begin{equation}
\bm{y}=\mathop{\Vert}_{n=1}^{N_\mathcal{Y}}\big[\spk{k_n}\,\tim{\hat a_n}\,\bm{y}^{(n)}\,\tim{\hat b_n}\big],
\label{eq:serial}
\end{equation}
where $\Vert$ concatenates and $\hat a_n,\hat b_n$ are $a_n,b_n$ in steps of $\delta_t$. The speaker token conditions
every translated word on its speaker, so two concurrent contributions become two segments; the time tokens give each
segment its source span. Targets use the diarizer's slots (each reference speaker is mapped to the slot most active
during their speech), so the Thinker learns to copy a slot rather than invent a label. The projector, the new
embeddings and LoRA adapters~\citep{hu2022lora} of the Thinker $p_\theta$ are trained with cross-entropy on $\bm{y}$,
with speaker and time tokens up-weighted because one wrong label misattributes a whole sentence.

\paragraph{Reading and writing.} The Thinker reads every $\Delta$ seconds, $t_r=\min(r\Delta,T)$. The diarizer finalises
a frame $\gamma$ seconds after it, so at read $r$ the speaker stream holds only frames up to $t_r-\gamma$. The Thinker
continues its committed output greedily until it emits $\wait$ (or $\eos$ at the last read) and commits what precedes
it. It learns when to write from deadlines: each reference translation is split at punctuation into chunks, and chunk
$j$, whose source ends at $e_j$, is due at read $r(j)=\min\{r:t_r\ge e_j+\epsilon\}$, with interpreting lag
$\epsilon$; the target of read $r$ is the chunks due at $r$, serialised by Eq.~\eqref{eq:serial} and followed by
$\wait$. This \emph{deadline Thinker}, trained on both directions, is shared by \sys{} and its baseline B, so the two
commit the same text with the same speakers at the same times, and all their differences arise afterwards.

\subsection{Speaker Streams and Playback Queues}
\label{sec:method-queue}

Committed text is heard only after the Talker has spoken it, and the blocks of one speaker are spoken one after
another. When a stream receives more speech than the source time it covers, later blocks wait; in B, two thirds of the
time from the end of a source unit until its translation is heard accrues after commitment
(Figure~\ref{fig:motivation}). We model this queue and let the speaking rate follow it.

\paragraph{Voice-prompted streams.} The text a slot commits at one read forms a \emph{block}; blocks of one speaker
committed in the same read are merged (\emph{read-level merging}). The UniSS Talker~\citep{uniss2025} speaks each block
as discrete speech codes, conditioned on its text, the stream's earlier speech, a speaking-rate token and a voice
prompt $\bm{v}_k$ cut from regions where slot $k$ alone is active, heard before the slot's first commitment
(Appendix~\ref{app:train-talker}).

\paragraph{Queue and heard delay.} In the stream $\mathcal{U}$ of one speaker, block $i=1,\dots,I$ holds $W_i$ target
units and is released at time $\nu_i$; a rendering $c$ of it lasts $D(c)$ seconds. The block waits a \emph{backlog}
$h_i$ for the previous block to finish, with $h_1=0$ and
\begin{equation}
h_{i+1}=\max\big(0,\;\nu_i+h_i+D(c_i)-\nu_{i+1}\big).
\label{eq:backlog}
\end{equation}
Spreading its units evenly over its playback, the delay summed over the block's units is
\begin{equation}
L_i(h_i,c_i)=W_i\,h_i+\tfrac{W_i+1}{2}\,D(c_i)
\label{eq:lat}
\end{equation}
unit-seconds. A long rendering thus costs twice: its own units are spread over more time, and through
Eq.~\eqref{eq:backlog} every later block waits longer.

\paragraph{Backlog-aware speaking rate.} B's base Talker uses one fixed rate, so a stream whose blocks arrive faster
than they are spoken builds a queue that only a pause in the source drains. The playback layer starts from the block's
\emph{pace} $q^{\mathrm{pace}}_i$, the rate at which its text would last as long as the source time since the previous
release, and raises it with the backlog:
\begin{equation}
q_i=q^{\mathrm{pace}}_i+q_0+\min\big(\bar q,\;\lfloor\kappa_h h_i\rceil\big),
\label{eq:rate}
\end{equation}
with offset $q_0$, $\kappa_h$ rate steps per second of backlog and a cap $\bar q$ that keeps speech intelligible. An
empty queue follows the source's pace; a stream that falls behind speaks faster until its queue drains. The rule,
however, reacts only to the existing backlog: it cannot see that a long rendering chosen now delays every later block,
a cost paid only in the stream's future.

\subsection{Playback Alignment: A Latency-Constrained CRF}
\label{sec:method-plem}

To charge each rendering for the delay it imposes on later speech, playback alignment (PA) scores whole playback
plans of a stream. Since a rendering affects later blocks only through the backlog it leaves, these plans form a chain,
and a linear-chain CRF~\citep{lafferty2001crf} over them admits exact inference.

\paragraph{Playback plans as a CRF.} For each block, the base Talker proposes a finite set $C_i$ of renderings, sampled at
the base rate and at a faster rate. A plan assigns block $i$ a state $z_i=(h_i,c_i)$; rounding Eq.~\eqref{eq:backlog} up
to a grid $\mathcal{H}=\{0,\delta_h,\dots,\bar H\}$ defines transitions $h_{i+1}=g_i(h_i,c_i)\in\mathcal{H}$, and a plan
$\bm{z}$ is feasible, $\bm{z}\in\mathcal{Z}(\mathcal{U})$, if $h_1=0$ and it follows $g_i$. Each state is scored by
\begin{equation}
\phi^\lambda_i(h,c)=\underbrace{\eta_\ell\,\bar\ell(c)+\eta_f\,f(c)}_{\text{quality}}-\underbrace{\lambda\,L_i(h,c)}_{\text{delay}},
\label{eq:plpot}
\end{equation}
where $\bar\ell(c)$ is the base Talker's mean log-probability per code, $f(c)$ the negative CTC loss per character of
the block's text under an independent recogniser~\citep{baevski2020wav2vec}, i.e.\ intelligibility, $\eta_\ell,\eta_f>0$
fixed weights, and $\lambda\ge0$ the price of a unit-second of delay. Normalising over feasible plans gives
\begin{equation}
p_\lambda(\bm{z}\mid\mathcal{U})=\frac{\mathds{1}[\bm{z}\in\mathcal{Z}(\mathcal{U})]}{Z_\lambda(\mathcal{U})}\exp\sum_{i=1}^{I}\phi^\lambda_i(z_i),
\label{eq:crf}
\end{equation}
a linear-chain CRF whose node potentials are Eq.~\eqref{eq:plpot} and whose transition potentials, $0$ or $-\infty$, encode
the stream's first-in-first-out queue.

\paragraph{Exact inference.} Sums over the $\prod_i|C_i|$ candidate sequences can therefore be grouped by the
backlog. The forward message $\alpha_i(h)$ (total weight of prefixes
leaving block $i$ with backlog $h$) and the backward message $\beta_i(h)$ (total weight of completions from there)
satisfy, with $\alpha_1(h)=\mathds{1}[h=0]$ and $\beta_{I+1}\equiv1$,
\begin{equation}
\begin{aligned}
\alpha_{i+1}(h')&=\textstyle\sum_{(h,c):\,g_i(h,c)=h'}\alpha_i(h)\,e^{\phi_i(h,c)},\\
\beta_i(h)&=\textstyle\sum_{c\in C_i}e^{\phi_i(h,c)}\,\beta_{i+1}\big(g_i(h,c)\big),\\
\mu_i(h,c)&=\alpha_i(h)\,e^{\phi_i(h,c)}\,\beta_{i+1}\big(g_i(h,c)\big)/Z_\lambda,
\end{aligned}
\label{eq:fb}
\end{equation}
where $Z_\lambda=\beta_1(0)$ and $\mu_i$ is the posterior marginal of state $(h,c)$; the expected delay of the stream
is $\mathbb{E}_{p_\lambda}[\Phi]=\sum_{i,h,c}\mu_i(h,c)L_i(h,c)$ with $\Phi(\bm{z})=\sum_iL_i(z_i)$. Inference is exact
within the state space and costs $O(I\,|\mathcal{H}|\max_i|C_i|)$ per stream (Proposition~\ref{prop:fb},
Figure~\ref{fig:pa}).

\paragraph{A latency budget instead of a tuned price.} The scale of $\lambda$ has no direct meaning, so we fix the delay
we accept and let $\lambda$ follow. Over a training corpus $\mathcal{S}$ of streams with $W$ target units, we seek the
distribution $q$ over plans closest to the quality-only CRF $p_0$ whose expected delay per unit stays within a budget
of $b$ seconds~\citep{ganchev2010posterior}:
\begin{equation}
\min_q\ \mathrm{KL}(q\,\Vert\,p_0)\quad\text{s.t.}\quad\mathbb{E}_q[\Phi_\mathcal{S}]\le bW,
\label{eq:pr}
\end{equation}
with $\Phi_\mathcal{S}$ the summed delay of all streams. As $\Phi_\mathcal{S}$ is a sum of per-block delays, the
solution is again the CRF of Eq.~\eqref{eq:crf} with one shared price $\lambda^\star$, the Lagrange multiplier of the
budget. The expected delay falls monotonically in $\lambda$, so bisection with Eq.~\eqref{eq:fb} finds $\lambda^\star$
(Proposition~\ref{prop:pr}). With $b$ set below the delay of base-rate playback, each direction obtains its own price
without tuning. The posterior then tells which renderings balance quality against present and future delay, but it
needs the release times of future blocks.

\subsection{Training the Talker on the CRF Posterior}
\label{sec:method-train}

Future release times are unknown during a meeting, so the CRF cannot run at test time. We use its posterior only in
training, as the target that teaches the Talker which renderings the CRF prefers given the backlog it faces.

\paragraph{Candidate weights.} On a training stream, we condition the marginals on the backlog $\tilde h_i$ that block
$i$ meets under base-rate playback; $\alpha_i(\tilde h_i)$ cancels and the candidates receive
\begin{equation}
w_i(c)=e^{\phi_i(\tilde h_i,c)}\,\beta_{i+1}\big(g_i(\tilde h_i,c)\big)\big/\beta_i(\tilde h_i),\quad c\in C_i.
\label{eq:w}
\end{equation}
Through $\beta_{i+1}$, a candidate is weighted by the quality and delay of every later block it makes wait, so a
rendering that is slightly better on its own but leaves a long backlog loses weight to a shorter one.

\paragraph{Objective.} The PA Talker $\pi_\xi$, initialised from the base Talker $\pi_0$, also reads a token
$\langle\tilde h_i\rangle$ of the binned backlog and is fine-tuned on the weighted candidates:
\begin{equation}
\begin{aligned}
\mathcal{J}(\xi)=\textstyle\sum_{\mathcal{U},i}\Big[&-\textstyle\sum_{c\in C_i}w_i(c)\log\pi_\xi(c\mid\cdot)\\
&+\kappa_q\,\mathrm{H}\big(\rho_i,\pi_\xi(q\mid\cdot)\big)\\
&+\kappa_{\mathrm{KL}}\,\mathrm{KL}\big(\pi_0\,\Vert\,\pi_\xi\big)\Big],
\end{aligned}
\label{eq:patrain}
\end{equation}
where $\pi_\xi(c\mid\cdot)$ is the probability of $c$'s codes given the block's text, voice prompt, history and backlog
token; $\rho_i(q)=\sum_cw_i(c)\mathds{1}[q(c)=q]$ is the posterior over rate tokens ($q(c)$ the rate $c$ was sampled
at), fitted by the cross-entropy $\mathrm{H}$; the KL keeps the codes close to $\pi_0$; and $\kappa_q,\kappa_{\mathrm{KL}}$
are weights. Up to a constant, the first term equals $\mathrm{KL}(w_i\Vert\tilde\pi_\xi)-\log\pi_\xi(C_i)$, with
$\tilde\pi_\xi$ the Talker renormalised on the candidates (Proposition~\ref{prop:patrain}): it fits the Talker's
preferences among renderings to the CRF posterior, one M-step of posterior-regularised EM~\citep{ganchev2010posterior}.
The target thus comes from a posterior that sees future release times, not from a teacher network
(cf.~\citealp{hu2016harnessing}).

\paragraph{Decoding.} At test time the PA Talker reads the backlog token of its actual queue, samples a few
candidates per block at the rate of Eq.~\eqref{eq:rate}, and keeps the one the recogniser finds most intelligible
(highest $f$). Read-level merging, the backlog-aware rate and the PA Talker form the \emph{playback layer}; \sys{} is
the deadline Thinker with this layer.

\subsection{Speaker-Attributed Speech Metrics}
\label{sec:method-metrics}

Text BLEU and text-emission latency show neither whether a translation is spoken into the right speaker's stream nor
when it is heard. We therefore score the synthesised speech automatically, through its ASR transcripts and heard times,
with counterparts of cpWER~\citep{watanabe2020chime6} and tcpWER~\citep{neumann2023meeteval}.

\paragraph{cpASR-BLEU.} Per meeting, the transcript $\hat{\bm{h}}_k$ of each output stream and the references
$\bm{r}_\varrho$ of each reference speaker are concatenated in time order and matched one-to-one by the assignment
$\sigma$ maximising the summed chrF~\citep{popovic2015chrf}. cpASR-BLEU is corpus BLEU over the matched pairs
$(\hat{\bm{h}}_k,\bm{r}_{\sigma(k)})$ together with every unmatched stream paired with an empty reference and every
unmatched speaker paired with an empty hypothesis. A sentence spoken into the wrong stream thus costs twice: it is
missing from its speaker's pair and superfluous in another's.

\paragraph{tcpASR-BLEU and latency.} tcpASR-BLEU adds the time constraint of tcpWER: each transcript segment, placed
at its source time with a collar, joins the overlapping reference unit of its matched speaker or counts as an
insertion, so a correct translation given to the wrong speaker or time earns nothing; on text output this gives
tcpBLEU (Appendix~\ref{app:eval}). Delay is measured from heard times: speech LAAL~\citep{papi2022over}, charging every
word of a block the time the block starts playing, and EndOffset, from a source unit's end until its last translated
word is heard.
